\section{Design of the \name file exchange system}
\label{design}
This section presents the \name protocol in detail.
Our goal is to enable private and reliable file exchange between users participating in the system, by using their respective devices only.
When Alice and Bob, two users of our system, want to exchange a file $f$ through \name; they start by sharing some information out-of-band, e.g. from their smartphones.
More precisely, they exchange metadata on $f$, encryption keys, and agree upon a Predictive Onion Route (POR).
PORs are our adaptation of onion routes, allowing secret communication between two entities despite the unreliability of the nodes that relay their messages.
From there on, Alice and Bob add the POR information to the headers of their subsequent file exchange messages, such that the file transfer remains secret to anyone but them.

% Alice and Bob each has a set of devices containing their files, \eg a 
% smartphone, laptop and a desktop computer.
% Alice and Bob meet in a café carrying only their smartphones and Alice wants to 
% share a file, \(f\), with Bob.
% Alice's smartphone knows some metadata of \(f\) and that it resides on Alice's 
% desktop computer.
% She gives some metadata about \(f\) to Bob, at least the file size.
% Bob creates an onion route of low-availability devices with some subset of Bob's 
% devices as the destination.
% The layers of the route are chosen such that the probability of all devices 
% being offline at the same time is low.
% Bob gives this route to Alice.
% Alice's smartphone ensures that her desktop computer (containing the file \(f\)) 
% gets the route and the instruction to send \(f\) using this route.
% When Alice's desktop computer receives this instruction it will create the first 
% half of the route and concatenate it with Bob's route, then transfer \(f\) to 
% Bob.
% In this section we will present the details of the necessary algorithms.

% We have a set of users, \(U\), each user \(u\in U\) has a set of devices, 
% \(D_u\).
% We call \(D_u\) the device swarm of \(u\).
% Alice, \(a\in U\), has a device swarm, \(D_a\), and Bob, \(b\in U\) has a device 
% swarm, \(D_b\).
% We let \(D = \cup_{u\in U} D_u\) be the set of all devices, the global device 
% swarm.
% Without loss of generality, we will use Alice, with device swarm \(D_a\), and 
% Bob, with device swarm \(D_b\), as examples.

\begin{figure}[t]
\centering
\def\svgwidth{\columnwidth}
%\input{outline.pdf_tex}
\import{figures/}{outline.pdf_tex}
%\includegraphics[width=0.9\columnwidth]{figures/outline.pdf}
\caption{\label{fig:outline}Outline of \name's components. Users squads of devices are circled in gray. From top to bottom: a user's devices collaborate through the Squad Overlay; online devices globally share information about them to random peers; using the two above overlays allows for the creation or Predictive Onion Routes (PORs) between sets of devices; secure file exchange between users is enabled by PORs.}
\end{figure}

The privacy objectives of our system require encryption primitives that we first detail in section~\ref{BuildingBlocks}. The other constituents of \name are depicted in figure~\ref{fig:outline}.
On top, the devices of each user are circled in gray, and form a user's device squad. The squad online members communicate among themselves: they exchange files and messages information, and cooperate to predict the user's future behavior. 
This first protocol is called the \emph{Squad Overlay} and will be discussed in section~\ref{sec:squad_overlay}.
%On top we see that each user's devices (shown as clouds) communicate with one another through a dissemination protocol---coined the \emph{squad overlay}---that will be discussed in section~\ref{sec:squad_overlay}. 
Below, we see all devices connected in a random graph: each device gives information about itself to random nodes when it is online. 
This \emph{Global Overlay} will be covered in section~\ref{sec:global_overlay}.
The information brought by these overlays allows devices to securely and reliably exchange messages between users by creating PORs, thus enabling file exchange between users, as will be explained in section~\ref{sec:message_passing}.
%Given this message passing protocol, users can privately exchange files, as shown in the bottom of the figure. We will cover this ultimate goal in section~\ref{sec:file_exchange}.


\include*{building-blocks}
\include*{squad-overlay-design}
\include*{global-overlay-design}
\include*{message-passing-design}
%\include*{file-exchange-design}